\scanpage{26}
% L1–4
\uline{Ex 1} Soit $\unsure{\mathfrak X|\mathscr X}$ un ens ordonné. Soit $D_*$ l'ens des parties finies de $\mathfrak X$, $\mathrm{tot}^{\mathrm t}$ ordonnées par l'ordre \unsure{induit|induite} (\textquotedblleft drapeaux de $\mathfrak X$\textquotedblright)\,; on ordonne $D_*$ par inclusion.
% L5–7
Alors $D_0$ (en tant qu'ens) \deleted{\unsure{est}} s'identifie à $\mathfrak X$, il est \deleted{\unsure{ordon|ordonné}} donc muni d'une relation d'ordre, et \textcircled{a} et \textcircled{b} sont évidemment satisfaites.
% L8–10
D'où un ens. quasi-simplicial $D_*$, et on a des homéom. can
\[
|D_*|\simeq|\unsure{\mathscr D|D}_*|\simeq|\mathfrak X|
\]
% L11–13
En fait, les $(D_*,\prec)$ obtenus sont ceux pour les quels on a, en plus de \textcircled{a}, \textcircled{b}, \inserted{\textcircled{c}} la condition
\[
\text{\textcircled{d}}\quad\forall d\in D_*,\ \text{on a}\quad d=\mathrm{Sup}(I_d)
\]
% L14–15
i.e. l'ens ordonné $D_*$ est l'ens des simplexes \inserted{non vides} d'une \uline{maquette simpliciale}, \deleted{$D_*$ sur \unsure{où l'ens quasi-simplicial des}}
% L16–20
Donc il faut prendre les maquettes simpliciales avec une structure d'ordre sur l'ens $\mathfrak X=D_0$ des sommets, induisant un ordre total sur chaque simplexe, et \deleted{\unsure{tel que les|tel que}} \inserted{(parmi les parties finies non vides)} sur les simplexes seulement\ldots
% L21–24
[Notons aussi que \deleted{moyennant \textcircled{a} et \textcircled{b} équivaut à \unsure{c}} \inserted{\textcircled{a} $+$ \textcircled{d} (i.e. $D_*$ correspond à} maquette simpliciale) équivaut à \textcircled{a$'$} $+$ \deleted{\unsure{\textcircled{c}}} \textcircled{d}, avec
% L24
\textcircled{a$'$} Dans $D_*$ les sup \inserted{(finis)} majorés existent\ldots\ ]
% archive sheet number: 25
